\section{Dữ liệu, phân tích khám phá và tiền xử lý}
\label{sec:data}

\subsection{Nguồn dữ liệu, biến đầu vào và quy trình tiền xử lý}
Bộ dữ liệu là bảng xếp hạng tiêu thụ nhiên liệu và phát thải \COtwo{} xe hạng nhẹ năm 2014 của Bộ Tài nguyên Thiên nhiên Canada~\cite{nrcan2014fuel} (tệp \texttt{FuelConsumptionCo2.csv}~\cite{ibm2014fuelco2}), gồm \nCars{} quan sát thuộc \nMakes{} hãng và \nCols{} thuộc tính, không có giá trị khuyết (\texttt{NaN}). Trên toàn tập dữ liệu, biến mục tiêu \texttt{CO2EMISSIONS} dao động từ \coMin{} đến \coMax{} g/km (trung bình \coMean{}, trung vị \coMedian{}, ĐLC \coSd{} g/km); dung tích động cơ (\texttt{ENGINESIZE}) có trung bình \engMean{} $\pm$ \engSd{} L (từ \engMin{} đến \engMax{} L); số xi-lanh từ \cylMin{} đến \cylMax{} (trung vị \cylMedian{}); và trong \nGeared{} xe hộp số phân cấp cố định (ngoài \nNoGear{} xe vô cấp \texttt{AV}), số cấp trải từ \gearsMin{} đến \gearsMax{} (trung vị \gearsMedian{}). Bảng~\ref{tab:vars-fe} hợp nhất toàn bộ thuộc tính gốc cùng các bước tiền xử lý và tạo đặc trưng.

\begin{table}[htbp]
  \centering
  \caption{Thuộc tính gốc trong \texttt{FuelConsumptionCo2.csv}, xử lý tạo đặc trưng và lý do thống kê.}
  \label{tab:vars-fe}
  \footnotesize
  \begin{tabular}{@{}l >{\raggedright\arraybackslash}p{5.5cm} >{\raggedright\arraybackslash}p{6.5cm}@{}}
    \toprule
    Cột gốc & Cách xử lý / đặc trưng dẫn xuất & Lý do kỹ thuật và thống kê \\
    \midrule
    \texttt{MODELYEAR} & Loại bỏ & Cột hằng (tất cả đều là năm 2014) \\
    \texttt{FUELCONSUMPTION\_*} & Loại bỏ cả 4 cột tiêu thụ nhiên liệu & Rò rỉ mục tiêu: \COtwo{} quy đổi trực tiếp từ đây (Mục~\ref{sec:leakage}) \\
    \texttt{ENGINESIZE}, \texttt{CYLINDERS} & Giữ nguyên dạng số; tạo bình phương và tương tác trong M4, M8b & Đặc trưng động cơ chính; giữ cả hai dù VIF sát 10 để kiểm định riêng và đối chiếu Ridge \\
    \texttt{FUELTYPE} & Mã hóa 3 biến giả (\texttt{Z}, \texttt{D}, \texttt{E}; tham chiếu \texttt{X}) & Tránh bẫy đa cộng tuyến; so sánh trực tiếp với xăng thường \\
    \texttt{VEHICLECLASS} & Gom \nVehicleClasses{} hạng thành \texttt{BODY} (5 mức; tham chiếu xe nhỏ) & Giảm từ 15 xuống 4 biến giả, tránh phình sai số chuẩn ở hạng ít mẫu \\
    \texttt{TRANSMISSION} & Tách thành \texttt{TRANS\_TYPE} (tham chiếu \texttt{A}), \texttt{GEARS}, \texttt{NOGEAR} & Xử lý khuyết có cấu trúc cho \nNoGear{} xe vô cấp \texttt{AV} (\texttt{GEARS}=0, \texttt{NOGEAR}=1) \\
    \texttt{MAKE}, \texttt{MODEL} & Không đưa trực tiếp; trích cờ \texttt{HYBRID} (tên chứa \texttt{HYBRID}, \texttt{PRIUS}, \texttt{CR-Z}, Lexus \texttt{...h}) & Sửa \nHybridMissed{} xe hybrid bị bỏ sót khi lọc theo chuỗi đơn thuần, tránh sai số hệ thống \\
    \texttt{CO2EMISSIONS} & Biến mục tiêu $y$ (giữ thang gốc g/km) & Giữ diễn giải vật lý; kiểm chứng riêng thang log ở M4log \\
    \bottomrule
  \end{tabular}
\end{table}

Toàn bộ bước mã hóa \texttt{OneHotEncoder} và chuẩn hóa \texttt{StandardScaler} (dùng cho Ridge M5 và Bayesian M6) được đóng gói trong \texttt{Pipeline} của \texttt{scikit-learn}~\cite{pedregosa2011sklearn} và chỉ khớp trên tập huấn luyện sau khi hoán vị ngẫu nhiên (\texttt{seed = 42}, chia 80/20 thành \nTrain{} xe train và \nTest{} xe test), ngăn ngừa rò rỉ thông tin từ tập kiểm tra.

\subsection{Phân tích khám phá dữ liệu (EDA)}

\subsubsection{Phân phối mục tiêu và phát hiện rò rỉ dữ liệu}
\label{sec:leakage}
Trên thang gốc, \texttt{CO2EMISSIONS} lệch phải nhẹ (độ lệch \coSkew{}, Jarque--Bera $JB = \jbRaw$, $p \jbRawP$), còn trên thang $\ln(\text{CO}_2)$ phân phối đối xứng hơn (độ lệch \logSkew{}, $JB = \jbLog$, $p = \jbLogP$); tuy nhiên, vì giả thiết chuẩn áp dụng cho phần dư chứ không cho phân phối biên của $y$, nhóm giữ thang gốc và kiểm chứng biến đổi log ở mô hình M4log.

Kiểm tra rò rỉ dữ liệu cho thấy tỉ số \texttt{CO2EMISSIONS} / \texttt{FUELCONSUMPTION\_COMB} gần như là hằng số tuyệt đối trong từng nhiên liệu: \ratioX{} (xăng thường \texttt{X}), \ratioZ{} (cao cấp \texttt{Z}), \ratioD{} (diesel \texttt{D}) và \ratioE{} (ethanol E85 \texttt{E}), với hệ số biến thiên không vượt quá \ratioCvMax\%. Như Hình~\ref{fig:leak-eda}(a) minh họa, các điểm nằm thẳng hàng trên ba đường tuyến tính; đưa \texttt{FUELCONSUMPTION\_COMB} và tương tác nhiên liệu vào OLS đạt $R^2 = \leakRtwo$ và RMSE $= \leakRmse$ g/km trên tập test. Vì \COtwo{} thực chất được quy đổi từ mức tiêu thụ nhiên liệu đo sau khi chạy thử, cả bốn cột \texttt{FUELCONSUMPTION\_*} bắt buộc phải bị loại.

\begin{figure}[htbp]
  \centering
  % Fig 02: leakage through fuel consumption, and CO2 vs engine size by body group.
% Legends placed outside below the axes so no data points are obscured.
\begin{tikzpicture}
\begin{groupplot}[group style={group size=2 by 1, horizontal sep=1.7cm},
  rpt, width=0.47\textwidth, height=5.5cm]
\nextgroupplot[title={(a) Rò rỉ qua mức tiêu thụ nhiên liệu},
  xlabel={Tiêu thụ nhiên liệu kết hợp (L/100 km)}, ylabel={Phát thải \COtwo{} (g/km)},
  ymin=80, ymax=530,
  legend style={at={(0.5,-0.25)}, anchor=north, legend columns=2, column sep=4pt}]
  \addplot[dot, cfuel-D] table[x=fc, y=co2] {\figdata{fig02-leak-D}};
  \addlegendentry{Diesel ($\approx$ \ratioD)}
  \addplot[dot, cfuel-X] table[x=fc, y=co2] {\figdata{fig02-leak-X}};
  \addlegendentry{Xăng thường (\ratioX)}
  \addplot[dot, cfuel-Z] table[x=fc, y=co2] {\figdata{fig02-leak-Z}};
  \addlegendentry{Xăng cao cấp (\ratioZ)}
  \addplot[dot, cfuel-E] table[x=fc, y=co2] {\figdata{fig02-leak-E}};
  \addlegendentry{Ethanol E85 (\ratioE)}
  \node[infobox, anchor=north west] at (rel axis cs:0.03,0.97)
    {OLS tương tác nhiên liệu:\\ $R^2_{\text{test}} = \leakRtwo$\\ RMSE $= \leakRmse$ g/km};
\nextgroupplot[title={(b) \COtwo{} theo dung tích và nhóm thân xe},
  xlabel={Dung tích động cơ (L)}, ylabel={Phát thải \COtwo{} (g/km)},
  xmax=9.0, ymin=80, ymax=530,
  legend style={at={(0.5,-0.25)}, anchor=north, legend columns=3, column sep=3pt}]
  \addplot[dot, cbody-small] table[x=eng, y=co2] {\figdata{fig02-body-small}};
  \addlegendentry{Xe nhỏ}
  \addplot[dot, cbody-midlarge] table[x=eng, y=co2] {\figdata{fig02-body-midlarge}};
  \addlegendentry{Trung--lớn}
  \addplot[dot, cbody-suv] table[x=eng, y=co2] {\figdata{fig02-body-suv}};
  \addlegendentry{SUV}
  \addplot[dot, cbody-pickup] table[x=eng, y=co2] {\figdata{fig02-body-pickup}};
  \addlegendentry{Bán tải}
  \addplot[dot, cbody-van] table[x=eng, y=co2] {\figdata{fig02-body-van}};
  \addlegendentry{Van \& ĐD}
  \addplot[cink, dashed, thick, domain=\pvEngMin:\pvEngMax, samples=2]
    {\pvMOneIntercept + \pvMOneSlope*x};
  \addlegendentry{Hồi quy M1}
\end{groupplot}
\end{tikzpicture}

  \caption{(a) Quan hệ gần như tất định giữa \COtwo{} và tiêu thụ nhiên liệu kết hợp theo loại nhiên liệu (rò rỉ mục tiêu). (b) \COtwo{} theo dung tích động cơ phân màu theo 5 nhóm thân xe, kèm đường hồi quy đơn M1.}
  \label{fig:leak-eda}
\end{figure}

\subsubsection{Phân nhóm, tương quan và đa cộng tuyến}
Bảng~\ref{tab:eda-groups} và Hình~\ref{fig:corr-body} cho thấy phát thải trung bình tăng dần theo nhóm thân xe từ xe nhỏ đến van \& đặc dụng. Đáng chú ý, trên trung bình thô, xe chạy ethanol E85 có \COtwo{} cao nhất (\meanE{} so với \meanX{} g/km ở xăng thường), nhưng dung tích trung bình của nhóm E85 cũng lên tới \engE{} L so với \engX{} L ở xăng thường --- bằng chứng sớm cho biến gây nhiễu sẽ được kiểm định ở Mục~\ref{sec:simpson}. Về tương quan (Hình~\ref{fig:corr-body}(a)), \texttt{ENGINESIZE} ($r = \corrEngCo$) và \texttt{CYLINDERS} ($r = \corrCylCo$) tương quan mạnh với \COtwo{} và với nhau ($r = \corrEngCyl$). Trong M3 trên tập train, $\mathrm{VIF} = \vifEngine$ (dung tích) và $\vifCyl$ (xi-lanh), sát nhưng dưới ngưỡng 10 (các biến thiết kế độc lập còn lại đều dưới \vifMaxOther); nhóm giữ cả hai biến, diễn giải thận trọng và đối chiếu với Ridge (M5).

\begin{table}[htbp]
  \centering
  \caption{Phát thải \COtwo{} trung bình và dung tích động cơ trung bình theo nhiên liệu, thân xe và hệ truyền động.}
  \label{tab:eda-groups}
  \footnotesize
  \begin{tabular}{@{}l l r r r@{}}
\toprule
Yếu tố & Nhóm & $n$ & \COtwo{} TB (g/km) & Dung tích TB (L) \\
\midrule
\multirow{4}{*}{Nhiên liệu} & Xăng thường & 514 & 241,1 & 2,95 \\
 & Xăng cao cấp & 434 & 268,5 & 3,63 \\
 & Diesel & 27 & 227,8 & 2,53 \\
 & Ethanol E85 & 92 & 291,1 & 4,44 \\
\midrule
\multirow{5}{*}{Nhóm thân xe} & Xe nhỏ & 355 & 235,6 & 2,93 \\
 & Sedan/wagon trung-lớn & 306 & 237,6 & 3,21 \\
 & SUV & 264 & 269,8 & 3,43 \\
 & Bán tải & 74 & 305,1 & 4,46 \\
 & Van \& đặc dụng & 68 & 342,2 & 4,61 \\
\midrule
\multirow{2}{*}{Hệ truyền động} & Không hybrid & 1033 & 259,2 & 3,37 \\
 & Hybrid & 34 & 166,2 & 2,56 \\
\bottomrule
\end{tabular}

\end{table}

\begin{figure}[htbp]
  \centering
  % Fig 03: Pearson correlation matrix, and CO2 by body group (boxplots).
\pgfplotstableread{\figdata{fig03-box}}\boxtable
\begin{tikzpicture}
\begin{groupplot}[group style={group size=2 by 1, horizontal sep=2.9cm}, rpt, half]
\nextgroupplot[title={(a) Ma trận tương quan Pearson},
  width=0.42\textwidth, height=0.42\textwidth,
  axis on top, grid=none, enlargelimits=false,
  xmin=-0.5, xmax=5.5, ymin=-0.5, ymax=5.5, y dir=reverse,
  xtick={0,...,5}, ytick={0,...,5},
  xticklabels={\COtwo, Dung tích, Xi-lanh, Số cấp, CVT (0 cấp), Hybrid},
  yticklabels={\COtwo, Dung tích, Xi-lanh, Số cấp, CVT (0 cấp), Hybrid},
  x tick label style={rotate=35, anchor=north east, inner sep=1pt},
  colormap={rdbu}{rgb255=(33,102,172) rgb255=(247,247,247) rgb255=(178,24,43)},
  point meta min=-1, point meta max=1,
  colorbar, colorbar style={width=2.5mm, ytick={-1,-0.5,0,0.5,1},
    yticklabel style={font=\tiny, /pgf/number format/use comma}}]
  \addplot[matrix plot*, mesh/cols=6, point meta=explicit]
    table[x=x, y=y, meta=r] {\figdata{fig03-corr}};
  \addplot[only marks, mark=none, point meta=explicit symbolic,
    visualization depends on={\thisrow{r} \as \rv},
    nodes near coords={\pgfmathparse{abs(\rv)>0.55}%
      \ifdim\pgfmathresult pt>0.5pt\color{white}\else\color{cink}\fi\pgfplotspointmeta},
    nodes near coords align=center, every node near coord/.append style={font=\tiny}]
    table[x=x, y=y, meta=label] {\figdata{fig03-corr}};
\nextgroupplot[title={(b) Phân phối \COtwo{} theo nhóm thân xe}, ygrid,
  width=0.46\textwidth, height=0.42\textwidth,
  ylabel={Phát thải \COtwo{} (g/km)},
  xtick={1,...,5}, xmin=0.4, xmax=5.6,
  xticklabels={Xe nhỏ, Trung--lớn, SUV, Bán tải, Van \& ĐD},
  boxplot/draw direction=y, boxplot/box extend=0.5]
  \preparedbox{\boxtable}{0}{draw=cink, solid, fill=cbody-small, fill opacity=0.7}
  \preparedbox{\boxtable}{1}{draw=cink, solid, fill=cbody-midlarge, fill opacity=0.7}
  \preparedbox{\boxtable}{2}{draw=cink, solid, fill=cbody-suv, fill opacity=0.7}
  \preparedbox{\boxtable}{3}{draw=cink, solid, fill=cbody-pickup, fill opacity=0.7}
  \preparedbox{\boxtable}{4}{draw=cink, solid, fill=cbody-van, fill opacity=0.7}
  \addplot[only marks, mark=*, mark size=0.9pt, cmuted, fill opacity=0.6, draw opacity=0]
    table[x=pos, y=y] {\figdata{fig03-box-out}};
\end{groupplot}
\end{tikzpicture}

  \caption{(a) Ma trận tương quan Pearson giữa các biến số thiết kế và \COtwo{}. (b) Phân phối \COtwo{} theo 5 nhóm thân xe sau khi gom từ \nVehicleClasses{} hạng xe gốc.}
  \label{fig:corr-body}
\end{figure}

\subsection{Hiệu chỉnh cờ hybrid từ chẩn đoán phần dư và xử lý xe flex-fuel}
% TODO(review): kiểm tra lại phần trình bày trước/sau bước sửa nhãn hybrid và ghi nhận xe flex-fuel.
Khi gán cờ hybrid đơn thuần bằng điều kiện tên mẫu chứa chuỗi \texttt{HYBRID}, chỉ nhận diện được \nHybridOld{} xe. Tuy nhiên, biểu đồ phần dư của mô hình M3 (Hình~\ref{fig:hybrid-fix}(a)) làm lộ một cụm \nHybridMissed{} xe không có chữ \texttt{HYBRID} trong tên (các mẫu Lexus hậu tố \texttt{h}, Toyota Prius, Honda CR-Z) nằm tách biệt ở vùng phần dư âm lớn tương tự nhóm hybrid đã gắn nhãn. Sau khi bổ sung quy tắc nhận diện (tổng \nHybrid{} xe hybrid), hệ số ước lượng đổi từ \bHybridOld{} thành \bHybrid{} g/km với khoảng tin cậy HC3 gọn hơn và BIC tập train giảm \bicDropFlag{} điểm (từ \bicThreeHOld{} xuống \bicThreeH{}, Hình~\ref{fig:hybrid-fix}(b)).

\begin{figure}[htbp]
  \centering
  % Fig 04: hybrid flag before / after the fix (P8).
\begin{tikzpicture}
\begin{groupplot}[group style={group size=2 by 1, horizontal sep=1.55cm}, rpt, height=4.5cm]
\nextgroupplot[title={(a) Phần dư M3 làm lộ các xe hybrid bị sót}, width=0.50\textwidth,
  xlabel={\COtwo{} dự đoán bởi M3 (g/km)}, ylabel={Phần dư M3 (g/km)},
  ymax=190, legend pos=north east]
  \addplot[cink, dashed, forget plot] coordinates {(100,0) (450,0)};
  \addplot[dot, csubtle, fill opacity=0.45] table[x=fit, y=res] {\figdata{fig04-regular}};
  \addlegendentry{Xe thường}
  \addplot[only marks, mark=*, mark size=1.7pt, cteal, draw=white, line width=0.3pt]
    table[x=fit, y=res] {\figdata{fig04-named}};
  \addlegendentry{Hybrid có chữ \texttt{HYBRID}}
  \addplot[only marks, mark=diamond*, mark size=2.4pt, cterra, draw=white, line width=0.3pt]
    table[x=fit, y=res] {\figdata{fig04-missed}};
  \addlegendentry{Hybrid bị sót (Lexus \dots h, Prius, CR-Z)}
\nextgroupplot[title={(b) Hệ số hybrid (KTC 95\% HC3) và BIC}, width=0.39\textwidth, ygrid,
  xlabel={Ước lượng hệ số hybrid (g/km)},
  xmin=-68, xmax=6, ymin=-0.6, ymax=1.6, ytick={0,1}, ymajorgrids=false, xmajorgrids=true,
  yticklabels={{Cờ cũ\\($n = \nHybridOld$)}, {Cờ đã sửa\\($n = \nHybrid$)}},
  yticklabel style={align=center}]
  \draw[csubtle, dashed] (axis cs:0,-0.6) -- (axis cs:0,1.6);
  \draw[cblue, very thick, |-|] (axis cs:\pvHybOldLo,0) -- (axis cs:\pvHybOldHi,0);
  \draw[cblue, very thick, |-|] (axis cs:\pvHybNewLo,1) -- (axis cs:\pvHybNewHi,1);
  \fill[cnavy] (axis cs:\pvHybOld,0) circle (2.2pt);
  \fill[cterra] (axis cs:\pvHybNew,1) circle (2.4pt);
  \node[lbl, anchor=north west, text=cmuted, font=\tiny] at (axis cs:-66,-0.22)
    {\bHybridOld{} g/km \;|\; BIC = \bicThreeHOld};
  \node[lbl, anchor=south west, text=cterra, font=\tiny\bfseries] at (axis cs:-66,1.22)
    {\bHybrid{} g/km \;|\; BIC = \bicThreeH{} ($-$\bicDropFlag)};
\end{groupplot}
\end{tikzpicture}

  \caption{Đối chiếu trước và sau bước sửa nhãn hybrid. (a) Phần dư M3 làm lộ \nHybridMissed{} xe hybrid không chứa chữ ``HYBRID'' trong tên. (b) Tác động lên $\hat\beta_{\mathrm{hybrid}}$ (KTC 95\% HC3) và BIC.}
  \label{fig:hybrid-fix}
\end{figure}

Ngoài ra, bộ dữ liệu chứa \nPairs{} mẫu xe flex-fuel được chứng nhận hai lần trên cả xăng thường (\texttt{X}) và ethanol E85 (\texttt{E}). Vì nhiên liệu là thông số đầu vào, nhóm giữ nguyên cả hai dòng, ghi nhận sự vi phạm một phần giả thiết độc lập (G4) ở phần hạn chế và sử dụng chính \nPairs{} cặp xe này cho kiểm định $t$ ghép cặp độc lập ở Mục~\ref{sec:simpson}.
